%% =========================================================
%% Chapter: The conditional movement estimate
%% Source: OpenAI, "A two-dimensional area law from a global spectral gap"
%%         (September 24, 2026), Lemma 5.1, 04-conditional.tex lines 14-308
%% =========================================================

\chapter{The Conditional Movement Estimate}
\label{ch:conditional_movement}

This chapter states the first conditional operator estimate
of~\cite[Lemma~5.1]{OpenAI2026AreaLaw}. A small subsystem $x$ of a pure state
changes sides between two systems $U$ and $P$; the norm of a vector obtained
by conjugating the state with marginal powers then decays at the rate of the
conditional information cost of the move. The constants are universal, and
the admissible exponents depend only on the dimension of $x$.

For a positive semidefinite matrix $\rho$ write $\Pi_{\ker\rho}$ for the
spectral projection onto its kernel and $\widehat\rho=\rho+\Pi_{\ker\rho}$ for
its kernel completion.

\begin{definition}[Kernel projection and kernel completion]
    \label{def:kernel_completion}
    \lean{Matrix.kernelProjection, Matrix.kernelCompletion}
    \leanok
    For a square complex matrix $\rho$, the kernel projection is
    $\Pi_{\ker\rho}=\mathbf 1_{\{0\}}(\rho)$, computed through the continuous
    functional calculus, and the kernel completion is
    $\widehat\rho=\rho+\Pi_{\ker\rho}$.
\end{definition}

\begin{lemma}[Real powers of the kernel completion]
    \label{lem:kernel_completion_rpow}
    \lean{Matrix.kernelProjection_mul_self,
        Matrix.kernelProjection_mul_kernelProjection,
        Matrix.cfc_rpow_kernelCompletion}
    \leanok
    \uses{def:kernel_completion}
    Let $\rho$ be Hermitian. Then $\Pi_{\ker\rho}\rho=0$,
    $\Pi_{\ker\rho}^2=\Pi_{\ker\rho}$, and for every real $s\ne0$,
    $\widehat\rho^{\,s}=\rho^s+\Pi_{\ker\rho}$, where $\rho^s$ vanishes on the
    kernel.
\end{lemma}

\begin{proof}\leanok
    All three identities are identities of functions on the spectrum of
    $\rho$; for the last one use $0^s=0$ when $s\ne0$.
\end{proof}

\begin{theorem}[Conditional movement estimate]
    \label{thm:conditional_movement}
    \lean{Entropy.conditionalMovement_norm_le}
    \leanok
    \uses{lem:kernel_completion_rpow, def:conditional_entropy}
    Let $\theta$ be a unit vector on four finite-dimensional systems
    $P,x,U,F$, and let $\rho_D$ denote its marginal on $D$. Set $d=\dim x$,
    $\ell=\log(ed)$ and $\eta=S(x|P)_\theta+S(x|U)_\theta$. There are
    universal constants $c,C>0$ such that, for $0<a\ell\le c$ and arbitrary
    positive semidefinite $\sigma$ on $xP$ and $\tau$ on $xU$ of trace at
    most one,
    \[
      \bigl\|\sigma^{a/2}\widehat\rho_P^{\,-a/2}
        \tau^{a/2}\widehat\rho_U^{\,-a/2}\theta\bigr\|
      \le\exp\Bigl(-\frac a2\eta+Ca^{5/4}\ell^2\Bigr).
    \]
    One may take $C=200000$ and $c=1/8$.
\end{theorem}

\begin{proof}\leanok
    Identify $\theta$ with its coefficient matrix from $xU$ to $PF$. The
    kernel projection of $\rho_U$ annihilates $\theta$, since its expectation
    in $\theta$ is $\operatorname{tr}(\Pi_{\ker\rho_U}\rho_U)=0$. The kernel
    projection of $\rho_P$ annihilates every vector obtained from $\theta$ by
    an operator on $xU$, because the two operators act on different factors.
    By the preceding lemma both kernel completions may therefore be replaced
    by the powers vanishing on the kernels. For these powers the estimate is
    the interpolation argument of~\cite[proof of Lemma~5.1]{OpenAI2026AreaLaw}:
    the analytic family
    $G(z)=\sigma^{az}\rho_P^{-az}\tau^{a(1-z)}\rho_U^{-a(1-z)}\theta$ on the
    strip $0\le\operatorname{Re}z\le1$ is bounded on the left boundary by the
    conditional Schatten estimate with the slack $\lambda(u)$ retained, the
    slack controls the movement of the state arriving at the right boundary,
    continuity of the conditional entropy with controlled system $x$ bounds
    the right boundary, and the three-lines inequality with harmonic weights
    $du/(2\cosh\pi u)$ combines the two. Faithful $\sigma,\tau$ are removed
    by regularization of their eigenvalues.
\end{proof}
